\subsection{Choosing the right collection (E3)}
\label{sec:routing-audit}
\label{sec:exp-e3}
The native deployment in E2 did not exercise catalogue-first routing.
Figure~\ref{fig:tiered-case} instead illustrates a separately captured
execution on BQ07. The agent used the catalogue and shelf descriptions to
select stored parts for the evidence filter, then read a bounded window of
the first part. The trace shows the steps the agent took, not how often this
path succeeds or whether it improves answer quality. It was not included in
the E2 judging or the E3 routing counts.

To assess collection selection, we ran a scripted selector with two
catalogue-card variants. A0 supplied the collection name and a description
derived from its contents; A3 was configured to supply the name alone. This
was not a clean description ablation: A3 could fall back to descriptions or
search outside the evaluated collections, and those paths were not logged.
The selector chose the labelled collection on all 72 questions with A0,
versus 37 with A3 in the frozen run (Figure~\ref{fig:routing}). On the same 72 questions, A0 remained at 72 while A3 reached 49 in each rerun.
The routing difference recurred, but the design does not isolate the effect of
descriptions.

Reaching the right collection is only a first step. Even with perfect A0
routing, the target string appeared in returned text on only 37 of 72
questions; correct collection selection alone did not ensure that the evidence
pack contained it.

\begin{figure}[!t]
\centering
\begin{minipage}{\columnwidth}
\centering
\includegraphics[width=0.96\linewidth]{fig_opt_route}
\caption{Correct collection selections out of 72 for A0 (name and description)
and A3 (nominally name-only) in the frozen, same-question follow-up and
later rerun. Bars give the number of questions in which the selector
chose the labelled collection; the denominator is fixed at 72 in every run.
A0 remains at 72/72, whereas A3 rises from 37/72 to 49/72 in both reruns.
Because fallbacks to descriptions or searches outside the evaluated collections
were not logged, this comparison is a routing diagnostic rather than a causal
description ablation.}
\label{fig:routing}
\end{minipage}
\end{figure}

\begin{figure*}[!t]
\centering
\includegraphics[width=\textwidth]{fig_opt_content}
\caption{Target-string containment on 72 in-domain questions. (a) Filled
circles and Wilson 95\% intervals show the frozen run; open triangles show
the same-question follow-up, with the adjacent columns giving the
corresponding hit counts. The dashed line marks frozen A0. A3 is an
auxiliary routing control, and read-all is a separate capped corpus scan
whose follow-up was incomplete. (b) For each displayed frozen-run contrast,
a filled circle counts questions hit only by the first configuration and an
open diamond counts questions hit only by the second; the connecting segment
connects the two discordant counts, and the symbols identify which
configuration each count favours. McNemar
$p$ values are exact, two-sided and unadjusted; values below 0.001 are
displayed as $<0.001$.}
\label{fig:content}
\end{figure*}

\subsection{Getting the target text into the evidence pack (E4)}
\label{sec:ladder-opt}
\label{sec:exp-e4}
\begin{table*}[!t]
\centering
\footnotesize
\setlength{\tabcolsep}{4pt}
\renewcommand{\arraystretch}{1.05}
\caption{Six main E4 configurations. A3 and read-all are separate
comparators. Limits are in characters, not tokens.}
\label{tab:ablation-configs}
\begin{tabular}{@{}l>{\raggedright\arraybackslash}p{0.15\textwidth}>{\raggedright\arraybackslash}p{0.77\textwidth}@{}}
\toprule
Arm & Search scope & Candidate selection and evidence\\
\midrule
A0 & One collection, selected by name and description & Retrieve vector hits. The reference-mode gate reads up to six distinct documents, with a nominal 40{,}000-character window per document.\\
A0f & Same route and initial search policy as A0 & Apply an inline gate to the candidates returned by that policy. Allow up to 250{,}000 characters per candidate and in the gate pack.\\
A0h & Same route as A0 & Keep the A0f gate and add a second vector search. Place up to eight chunk windows first, then add gate-retained text up to a final limit of 60{,}000 characters.\\
PD & All ten collections & Keep the six highest-scoring distinct collection/path pairs and use the A0f inline gate.\\
A1b & All ten collections & Merge vector hits and use the A0 reference-mode gate, including document reads.\\
A2 & All ten collections & Merge the same hits as A1b and concatenate their snippets. Do not read documents or run a gate.\\
\bottomrule
\end{tabular}
\end{table*}

We next varied search scope and evidence assembly. A3 uses A0's downstream
evidence construction but changes the route cards and collection setup. It is an auxiliary routing diagnostic; the mechanical
corpus scan is a separate comparator. The six main configurations differ
in search scope (one collection selected by name and description, or all
ten), in whether documents are read and gated, and in how the pack is
ordered. Routed vector search requests up to 30 candidates at a cosine
threshold of 0.35 and retries at 0.30 only if it returns nothing; the
unrouted arms search each collection at 0.35. A0, A1b and A3 could use
the gold domain to resolve document references, but not to choose a route card
(Section~\ref{sec:threats}).

A0h had the highest target-string count among the six main evidence-pack
configurations in the frozen run: 52 of 72 questions, up from 37 with A0.
It found the target string on 15 questions where A0 missed it and lost no
A0 hit (exact McNemar $p=6.1\times10^{-5}$). Both reruns on the same questions retained this advantage: 51 versus 39 in
the follow-up and 48 versus 36 in the later rerun (Supplementary Table~S12).
A0h also led the unrouted vector union A1b. Its difference from the unrouted
per-document configuration PD was not significant ($p=0.070$).

These contrasts compare configurations rather than single components. The
route and initial search policy are shared by A0, A0f and A0h, but the
realised candidate lists can still differ. A0h changes both what the system
retrieves and what it puts first in the pack, so its comparison with A0f is
not a test of window size alone. A2 removes
document reads as well as the gate, so it does not isolate the gate's
effect. The arm-specific trace counters were 309~kB per question for A0h, 31.4~kB
for A0 and 24.3~MB for the capped corpus scan, which found every target.
These figures use different logging boundaries and are not comparable measures
of I/O or generator input. The scan capped each document read at
60{,}000 characters. A0h's counter sums fetched candidate text, gate output
and selected windows, so it can count text twice; A0's counter excludes the
gate's internal reference-mode reads. A0h's final pack has a separate
60{,}000-character cap. Table~\ref{tab:ablation-configs} gives the limits by
configuration; Figure~\ref{fig:content} shows the containment contrasts.

\begin{figure*}[t]
\centering
\includegraphics[width=\textwidth]{fig_opt_verifier}
\caption{Ordering by verifier score lost eight top-five hits and gained none
in a replay of 16 fixed candidate pools. (a) Empirical cumulative distributions
of the 191 recorded Laya scores (17 gold and 174 non-gold; medians 0.8759 and
0.86265, respectively). All recorded scores exceeded the 0.5 cutoff; Q14
had one unrecorded score. (b) Each matrix cell gives the count and IDs of
pools with the indicated combination of hit or miss under the two orders.
A hit means that at least one gold document was in the top five. Fusion hit
in 15/16 pools and verifier-score order in 7/16 (exact McNemar $p=0.0078$).
The replay tests ordering, not passage admission.}
\label{fig:verifier}
\end{figure*}

\subsection{Storing and finding operational lessons (E5)}
\label{sec:e3}
\label{sec:exp-e5}
HMRS stores operational lessons as searchable documents and orders them by
credibility tier. Only parts of this workflow could be verified. A smoke test
produced candidate drafts but did not establish global recall. A separate
case report placed a stored lesson first in a search; no answer trace was
retained. The
approval check found an empty draft pool, and the heuristic dry run did not
meet its quality criteria. These checks illustrate parts of the workflow;
they do not show reliable recovery or automatic approval of lessons in
deployment.
